\section{Benchmark Case Studies：从目标到数据生成与验证}

课程没有只列 benchmark 名称，而是先给出构造问题：目标是什么、Agent 做什么、benchmark 如何评价、数据怎样生成、如何保证质量。读案例时应比较其 task pipeline、oracle、真实性和污染风险，而不是只看榜单。

\subsection{案例方法：先画数据与执行流水线}

slides 48 与 50 把 benchmark 设计拆成目标、任务、测试过程和 data pipeline。一个可维护 benchmark 需要明确任务来源、环境初始化、participant interface、结果收集、评分和质量审查。若数据生成不可复现，分数波动无法归因；若任务目标模糊，再精确的评分也没有意义。

\lecturefigure{slides-images/slide-048.png}{官方 slide 48：案例分析先问 goal、task 与 evaluation}
\lecturefigure{slides-images/slide-050.png}{官方 slide 50：good benchmark 的构造与 data pipeline}

\begin{importantbox}{Benchmark datasheet 最小字段}
任务意图、来源与许可；环境与依赖版本；Agent 可见信息和工具；oracle/metric；成功与失败样例；难度与分布；污染风险；成本预算；已知限制。把这些字段写成机器可读 metadata，才能支持重跑、切分和跨版本比较。
\end{importantbox}

\subsection{CyberGym：漏洞复现中的容器、补丁与 verifier}

slides 51--53 展示 CyberGym：从真实漏洞/补丁构造容器化任务，让 Agent 在代码环境中发现或利用漏洞，并用 PoC、patch 或测试验证。其价值是任务与真实安全问题相关；难点是构建环境、依赖、漏洞可复现性和安全隔离。补丁本身可能泄漏答案，数据生成必须防止把 ground truth 暴露给 participant。

\lecturefigure{slides-images/slide-051.png}{官方 slide 51：CyberGym task、environment、PoC 与 patch pipeline}
\lecturefigure{slides-images/slide-052.png}{官方 slide 52：CyberGym 模型表现与案例结果}
\lecturefigure{slides-images/slide-053.png}{官方 slide 53：CyberGym 数据生成、复现与 contamination analysis}

\begin{knowledgebox}{读图：安全 benchmark 的双重边界}
一方面要让漏洞真实可执行，另一方面不能让实验逃逸 sandbox 或泄漏危险能力。环境应禁用外网、限制资源、记录系统调用并销毁实例；评估应区分找到漏洞、生成有效 exploit、修复漏洞和不破坏其他功能。单一成功率不能概括不同安全目标。
\end{knowledgebox}

\subsection{tau-bench 与 tau2-bench：对话、工具和双控制}

slides 54--57 描述零售/航空等场景中的对话 Agent。tau-bench 用用户模拟、policy、数据库/API 和任务结果评估工具使用；tau2-bench 进一步让用户侧也能调用工具，形成 dual-control environment，更接近双方都可改变状态的真实流程。数据由 schema、policy 和任务模板生成，成功由最终数据库状态与约束共同判断。

\lecturefigure{slides-images/slide-054.png}{官方 slide 54：tau-bench 的用户、Agent、工具与数据库}
\lecturefigure{slides-images/slide-055.png}{官方 slide 55：tau-bench 数据生成、tool schema 与 pass@k}
\lecturefigure{slides-images/slide-056.png}{官方 slide 56：tau2-bench 双控制交互结构}
\lecturefigure{slides-images/slide-057.png}{官方 slide 57：tau2-bench 数据、验证和可靠性度量}

\begin{knowledgebox}{Pass@k 与可靠性}
若同一任务运行 $k$ 次，pass@k 常衡量至少一次成功的概率；生产系统还关心 pass$^k$ 或多次都成功的稳定性。Agent 有随机策略和模拟用户，两者都会贡献方差。报告平均成功率时应同时给重复运行分布、失败类型和成本，否则“一次能做成”会被误解为可靠服务。
\end{knowledgebox}

\begin{warningbox}{Simulator realism 是隐藏变量}
用户模拟器若过于合作，会高估 Agent；若语言模式单一，模型可识别 benchmark；policy 与数据库若简化，也会遗漏真实异常。tau 类 benchmark 支持受控比较，不证明所有客服流程都能自动化。应加入 adversarial user、工具故障、权限限制和真实对话回放。
\end{warningbox}

\subsection{GDPval：专业任务与真实经济价值}

前几个案例主要依赖程序化环境，GDPval 则把评估推向专业知识工作。本小节需要同时观察 deliverable 质量、领域专家判断和任务的经济相关性，并警惕把一份可评分产出直接外推为岗位级自动化结论。
slides 58--59 介绍 GDPval，任务来自不同职业的知识工作，强调 deliverable 的现实性和专业评审。它试图把模型输出与经济活动联系起来，而非只测学术题。难点在于专家时间昂贵、任务可能包含机密背景、rubric 主观且不同职业难以放在同一标尺。

\lecturefigure{slides-images/slide-058.png}{官方 slide 58：GDPval 的职业任务与经济相关性}
\lecturefigure{slides-images/slide-059.png}{官方 slide 59：GDPval 数据生成、专家评审与局限}

\begin{knowledgebox}{读图：Economic relevance 不等于自动化比例}
benchmark 能测某些专业 deliverable 的质量，不直接告诉我们岗位有多少工作可被替代。真实生产还包含沟通、责任、隐性知识、数据访问和组织流程。GDPval 更适合比较模型在工作样本上的能力边界，并发现哪些任务需要专家监督。
\end{knowledgebox}

\subsection{CRMArena：业务流程与 hidden state}

slides 60--61 将 Agent 放入 CRM 场景，要求理解客户、销售和服务数据并执行多步操作。环境状态复杂、字段关系多，成功不仅是语言正确，还要数据库变更符合业务规则。数据 pipeline 需要 schema、任务模板、初始状态和 ground truth transition。

\lecturefigure{slides-images/slide-060.png}{官方 slide 60：CRMArena 的 CRM 环境、角色与任务}
\lecturefigure{slides-images/slide-061.png}{官方 slide 61：CRMArena 数据生成、验证与限制}

\begin{warningbox}{业务 benchmark 的权限与幂等性}
真实 CRM 有角色权限、审批、重复提交和审计要求。benchmark 若允许 Agent 直接写所有字段，会高估能力并忽略风险。应检查动作是否授权、重复执行是否安全、无关记录是否被修改，以及失败后能否回滚。
\end{warningbox}

\subsection{LegalAgentBench：长文档、检索与专业判断}

slides 62--63 介绍法律 Agent 任务，涉及检索、引用、问题回答和程序化流程。法律文本长、术语精确、时效性强，数据需要律师或可靠来源验证。字符串匹配难以评价论证，LLM judge 又可能被流畅度欺骗，因此要组合 citation validity、source coverage、规则检查和专家复核。

\lecturefigure{slides-images/slide-062.png}{官方 slide 62：LegalAgentBench 的法律任务与工具环境}
\lecturefigure{slides-images/slide-063.png}{官方 slide 63：LegalAgentBench 数据生成、验证与限制}

\begin{knowledgebox}{案例对照：oracle 从哪里来}
\begin{tabular}{p{2.5cm}p{3.9cm}p{4.7cm}}
\toprule
\textbf{Benchmark} & \textbf{主要 oracle} & \textbf{最难边界} \\
\midrule
CyberGym & PoC、patch、tests & 隔离、复现、答案泄漏 \\
tau/tau2 & final DB state、policy & 用户模拟与多次可靠性 \\
GDPval & 专家 rubric & 主观性、成本、职业可比性 \\
CRMArena & state transition、规则 & 权限、幂等、隐藏关系 \\
Legal Agent & 引用、规则、专家 & 时效、长文档、责任 \\
\bottomrule
\end{tabular}
\end{knowledgebox}

\subsection{本章小结}

案例表明 benchmark 质量来自任务目标、环境、数据 pipeline 与 oracle 的协同。每个案例都牺牲部分真实性换取复现性，因此必须公开限制。下一章将 benchmark 包装成 Green Agent，使它可以在 AgentBeats 式平台上自动运行和汇报。

